% =====================================================================
%  BLOCCHI FIGURA per report_it.tex
%
%  Ognuno sostituisce un \figph{...} esistente, oppure si inserisce nel
%  punto indicato. Nessuna riga di prosa viene toccata.
%
%  I file immagine vanno caricati su Overleaf nella stessa cartella di
%  phantomx.png (la radice del progetto), con questi nomi:
%     rampa.pdf  ostacolo.pdf  curva.pdf  impulso.pdf  inviluppo.pdf
%     fattibilita.png
%  stanno in  grafici/report/  tranne fattibilita.png che sta in grafici/
% =====================================================================


% ---------------------------------------------------------------------
% [1]  SOSTITUISCE il \figph{70mm}{fig:obstacle}{...}
%      (sezione "Discontinuita': la ricerca del terreno ripaga")
% ---------------------------------------------------------------------
\begin{figure}[htbp]
  \centering
  \includegraphics[width=\linewidth]{ostacolo.pdf}
  \caption{Beccheggio del corpo durante il superamento dell'ostacolo singolo
  (T5). In grigio la finestra in cui almeno un piede si trova sull'ostacolo,
  identificata cinematicamente. Fuori da quella finestra le tre curve
  oscillano attorno allo zero con l'ampiezza del passo; dentro, \Co{} e \Ct{}
  compiono un'escursione di diversi gradi in entrambi i versi --- prima il
  muso in gi\`u mentre le zampe anteriori salgono sul gradino, poi in su quando
  vi salgono le posteriori --- mentre \Cth{} resta entro circa un grado per
  tutto il passaggio. \`E la stessa differenza che la Tabella~\ref{tab:t4dt5}
  riporta come $-73\%$ di beccheggio di picco.}
  \label{fig:obstacle}
\end{figure}


% ---------------------------------------------------------------------
% [2]  SOSTITUISCE il \figph{70mm}{fig:ramp}{...}
%      (sezione "La rampa")
% ---------------------------------------------------------------------
\begin{figure}[htbp]
  \centering
  \includegraphics[width=\linewidth]{rampa.pdf}
  \caption{Inclinazione del corpo durante la salita della rampa di
  \SI{8}{\degree} (T4). La linea tratteggiata a \SI{8}{\degree} \`e la pendenza
  della rampa: una curva che vi si appoggia descrive un tronco \emph{parallelo
  al terreno}. I primi \SI{4.5}{\second} sono in piano; poi \Co{} e \Ct{}
  salgono fino alla pendenza e vi restano per il resto della corsa, mentre
  \Cth{} si ferma sotto \SI{1.5}{\degree}, cio\`e tiene il tronco orizzontale.
  Non \`e un'inclinazione pi\`u piccola della stessa cosa: sono due comportamenti
  diversi, ed \`e il motivo per cui la riga corrispondente della
  Tabella~\ref{tab:t4} non porta una percentuale. Il dente di sega
  sovrapposto a tutte e tre le curve \`e l'oscillazione propria del passo.}
  \label{fig:ramp}
\end{figure}


% ---------------------------------------------------------------------
% [3]  SOSTITUISCE il \figph{65mm}{fig:envelope}{...}
%      (sezione "L'inviluppo di stabilita' in velocita'")
%      LEGGI LA NOTA NEL MESSAGGIO: i valori della figura non coincidono
%      con quelli della Tabella~\ref{tab:envelope}.
% ---------------------------------------------------------------------
\begin{figure}[htbp]
  \centering
  \includegraphics[width=\linewidth]{inviluppo.pdf}
  \caption{Spazzata in velocit\`a (T2) per i tre controllori: rimbalzo verticale
  del corpo (sopra) e frazione del task percorsa (sotto) contro il fattore di
  velocit\`a comandato. Il minimo del rimbalzo non sta agli estremi ma attorno al
  punto di lavoro nominale: a met\`a velocit\`a il corpo oscilla perch\'e il passo
  diventa lungo rispetto al tempo di appoggio, oltre il nominale perch\'e la
  zampa non fa in tempo a riportarsi sotto il corpo. La frazione del task
  segue: resta attorno al $100\%$ nelle celle centrali e crolla al raddoppio
  della velocit\`a, dove diventa negativa --- il robot arretra invece di
  avanzare.}
  \label{fig:envelope}
\end{figure}


% ---------------------------------------------------------------------
% [4]  SOSTITUISCE il \figph{65mm}{fig:feasibility}{...}
%      (sezione "Fattibilita' sugli attuatori reali")
% ---------------------------------------------------------------------
\begin{figure}[htbp]
  \centering
  \includegraphics[width=\linewidth]{fattibilita.png}
  \caption{Coppia richiesta ai giunti contro i limiti dell'AX-12A, per task e
  controllore. Vanno letti due regimi distinti: la coppia efficace del giunto
  pi\`u caricato, che descrive la marcia e resta stabilmente sotto lo stallo ma
  sopra il carico continuo raccomandato dal costruttore, e la coppia di picco,
  che lo stallo lo supera in corrispondenza degli urti all'appoggio. I valori
  numerici corrispondenti sono nella Tabella~\ref{tab:feasibility}.}
  \label{fig:feasibility}
\end{figure}


% ---------------------------------------------------------------------
% [5]  NUOVA. Inserire nella sezione "Piano e curva: la ricerca del terreno
%      costa assetto", DOPO il paragrafo che comincia con
%      "Su un terreno che e' esattamente quello che il generatore..."
%      (cioe' subito prima del \figph di fig:pitch-flat, oppure dopo)
% ---------------------------------------------------------------------
\begin{figure}[htbp]
  \centering
  \includegraphics[width=\linewidth]{curva.pdf}
  \caption{Traiettoria del centro di massa nel piano orizzontale con comando
  di imbardata $+\SI{0.1}{\radian\per\second}$ (T3). Le tre curve si
  sovrappongono quasi ovunque: il comando di imbardata viene seguito allo
  stesso modo dai tre controllori, e le differenze della
  Tabella~\ref{tab:t2t3} su questo task riguardano l'assetto del corpo, non il
  percorso. \textbf{Le scale dei due assi sono diverse} --- quasi due metri in
  ascissa contro quattro decimetri in ordinata --- quindi lo scostamento
  laterale appare molto pi\`u marcato di quanto sia.}
  \label{fig:curva-t3}
\end{figure}


% ---------------------------------------------------------------------
% [6]  NUOVA. Inserire nella sezione "Impulso laterale: una lacuna comune a
%      tutti e tre", DOPO il paragrafo "A impulso grande le tre curve
%      coincidono entro il 2%..."
% ---------------------------------------------------------------------
\begin{figure}[htbp]
  \centering
  \includegraphics[width=\linewidth]{impulso.pdf}
  \caption{Deviazione laterale massima contro l'ampiezza dell'impulso
  applicato (T7), in scala logaritmica sull'ascissa per coprire le tre decadi
  della spazzata. Le tre curve sono praticamente indistinguibili su tutto
  l'intervallo: fino a circa \SI{0.7}{\newton\second} la deviazione resta di
  pochi millimetri per tutti, oltre cresce bruscamente fino a quasi
  \SI{400}{\milli\meter} senza che nessuno dei tre controllori si distingua
  dagli altri. \`E la rappresentazione diretta di una lacuna architetturale
  comune: n\'e la ricerca del terreno n\'e l'anello d'assetto intervengono su
  questo disturbo.}
  \label{fig:impulso-t7}
\end{figure}

% ---------------------------------------------------------------------
% [7]  SOSTITUISCE il \figph{58mm}{fig:c3p-pacco}{...} in sezione_c3p.tex
%      File immagine: pacco.pdf  (sta in grafici/report/)
%      Richiede di rilanciare T4, T4D e T6 in C3P con la versione di
%      stato_pacco.m che salva la serie nel tempo.
% ---------------------------------------------------------------------
\begin{figure}[htbp]
  \centering
  \includegraphics[width=\linewidth]{pacco.pdf}
  \caption{Spostamento del pacco sul vassoio rispetto alla posizione di
  riferimento, per i tre task, con la soglia di \SI{50}{\milli\meter}
  corrispondente al gioco fra cubo e vassoio. Le tre curve restano per tutta
  la corsa nel quinto inferiore del riquadro: \`e la rappresentazione diretta
  del fattore due e mezzo di margine. La forma delle curve aggiunge
  un'informazione che la Tabella~\ref{tab:c3p} non pu\`o portare: su T4D il
  pacco raggiunge il massimo a circa \SI{17}{\second} e poi \emph{rientra},
  mentre su T4 e T6 il massimo cade negli ultimi istanti della corsa, cio\`e
  lo spostamento sta ancora crescendo quando la simulazione termina. Il
  margine misurato \`e quindi un margine \emph{alla fine della run}, non un
  valore di regime, e corse pi\`u lunghe andrebbero verificate.}
  \label{fig:c3p-pacco}
\end{figure}
